%%%%%%%%%%%%%%%%%%%%%%%%%%%%%%%%%%%%%%%%%%%%%%%%%%%%%%%%%%%%%%%%%%%%%%%%%%%%%%%%%%%%%%%%%%%%%%%%%%%%%%%%%%%%%%%%%%%%%%%%%%%%%%%%%%%%%%%
\cvsection{Professional Experience}

\cvevent{\href{https://isrp.pt/integrated-members/}{}}{Robotics Engineer - \href{https://www.oceanscan-mst.com/}{OceanScan-MST}}{September 2025 -- Ongoing}{Porto de Leixões, Porto}
\begin{itemize}
\item Development/testing/validation of AHRS calibration techniques for improved AUV navigation.
\item Development/testing/validation of algorithms for reliable situational awareness and safe/robust control of AUVs, under the \href{https://lsts.fe.up.pt/systems}{DUNE/Neptus} software toolchain.
% (C++/Java).
\end{itemize}

\cvevent{\href{https://sigarra.up.pt/feup/en/FUNC_GERAL.FORMVIEW?p_codigo=620202}{}}{Invited Auxiliar Professor - FEUP}{Feb. 2022 -- Ongoing}{University of Porto}
\begin{itemize}
	\item Linear Algebra: theory of matrices, solutions of linear systems of equations, vector spaces, linear transformations, eigenvalue problems.
	\item Computing Systems: fundamentals of concurrent computing. Assistance in laboratory practices with Arduino FreeRTOS and Linux.
\end{itemize}

\cvevent{\href{https://isrp.pt/integrated-members/}{}}{Researcher at \href{https://isrp.pt/}{ISR (Institute for Systems and Robotics)}}{Apr 2019 -- Aug 2025}{University of Porto}
\begin{itemize}
\item Research areas: formal methods for safe, deadlock-free control for cyber-physical systems and autonomous navigation techniques with collision avoidance guarantees for multi-robot systems.
\item Conducted original research in engineering and published papers in top journals (IFAC Automatica, IEEE Robotics and Automation Letters) and conferences in robotics (IROS, ICRA, among others).
\end{itemize}

\cvevent{\href{https://cordis.europa.eu/project/id/642153/reporting}{}}{Researcher at the MarineUAS Project - \href{https://systec.fe.up.pt/}{SYSTEC}}{Oct 2018 -- March 2019}{University of Porto}
\begin{itemize}
\item Research in nonlinear, cooperative and autonomous control for multi-robot systems.
\item \href{https://lsts.fe.up.pt/news/meet-matheus-f-reis}{Development and testing of robust moving path-following control strategies for teams of Light Autonomous Underwater Vehicles (LAUVs)}.
\end{itemize}

%%%%%%%%%%%%%%%%%%%%%%%%%%%%%%%%%%%%%%%%%%%%%%%%%%%%%%%%%%%%%%%%%%%%%%%%%%%%%%%%%%%%%%%%%%%%%%%%%%%%%%%%%%%%%%%%%%%%%%%%%%%%%%%%%%%%%%%
\cvsection{Research Projects}

\cvevent{Advances in Learning and Control Theory for Safety Critical Multi-Robot Systems}{\href{https://systec.fe.up.pt/}{SYSTEC}}{Jan 2021 -- Dec 2024}{Porto}
\begin{itemize}
\item The research has developed tools and algorithms for provably safe and deadlock-free operation of safety-critical systems 
using the unified framework of Control Lyapunov Functions (CLFs) and Control Barrier Funcions (CBFs). Position: PhD candidate.
\end{itemize}

% \divider
\newpage

\cvevent{STRIDE}{\href{https://systec.fe.up.pt/}{SYSTEC}}{April 2019 -- Dec 2020}{Porto}
\begin{itemize}
\item The STRIDE project focused in R\&D in Systems, Control and Optimization for Networked Robotic Vehicle Systems and Advanced Production Systems, Mobility, Energy and Environment.
Position: Early stage researcher.
\end{itemize}

% \newpage
% \divider

\cvevent{\href{https://cordis.europa.eu/project/id/642153/reporting}{MarineUAS - Autonomous Unmanned Aerial Systems for Marine and Coastal Monitoring}}{\href{https://systec.fe.up.pt/}{SYSTEC}}{Oct 2018 -- March 2019}{Porto}
\begin{itemize}
\item Innovative Training Network on Autonomous Unmanned Aerial Systems for Marine and Coastal Monitoring, funded by Marie Skłodowska-Curie Actions.
Position: Early stage researcher.
\end{itemize}

% \divider 

\cvevent{HEADS}{\href{https://www.embrapii.coppe.ufrj.br/laboratorios/1217/lead-laboratorio-de-controle-e-automacao-engenharia-de-aplicacao-e-desenvolvimento}{Coppetec/COPPE-UFRJ}}{Oct 2016 -- Aug 2018}{Rio de Janeiro}
\begin{itemize}
\item The project HEADS (Hydrocarbon Early Automatic Detection System) is a collaboration among several research centers with the objective of developing a 3 degree-of-freedom stabilization system for a thermal camera and test the integrated system on a marine vehicle. Position: M.Sc student.
\end{itemize}

\medskip

%%%%%%%%%%%%%%%%%%%%%%%%%%%%%%%%%%%%%%%%%%%%%%%%%%%%%%%%%%%%%%%%%%%%%%%%%%%%%%%%%%%%%%%%%%%%%%%%%%%%%%%%%%%%%%%%%%%%%%%%%%%%%%%%%%%%%%%
\cvsection{Publications}

%% Specify your last name(s) and first name(s) as given in the .bib to automatically bold your own name in the publications list.
%% One caveat: You need to write \bibnamedelima where there's a space in your name for this to work properly; or write \bibnamedelimi if you use initials in the .bib
%% You can specify multiple names, especially if you have changed your name or if you need to highlight multiple authors.
% \mynames{Lim/Lian\bibnamedelima Tze,
%   Wong/Lian\bibnamedelima Tze,
%   Lim/Tracy,
%   Lim/L.\bibnamedelimi T.}
%% MAKE SURE THERE IS NO SPACE AFTER THE FINAL NAME IN YOUR \mynames LIST

\nocite{*}

% \printbibliography[heading=pubtype,title={\printinfo{\faBook}{Books}},type=book]

% \divider

\printbibliography[heading=pubtype,title={\printinfo{\faFile*[regular]}{Journal Articles}},type=article]

% \divider

\printbibliography[heading=pubtype,title={\printinfo{\faUsers}{Conference Proceedings}},type=inproceedings]

% \divider

%%%%%%%%%%%%%%%%%%%%%%%%%%%%%%%%%%%%%%%%%%%%%%%%%%%%%%%%%%%%%%%%%%%%%%%%%%%%%%%%%%%%%%%%%%%%%%%%%%%%%%%%%%%%%%%%%%%%%%%%%%%%%%%%%%%%%%%
% \cvsection{Teaching Experience}

% \cvevent{Assistant Invited Professor at the Faculty of Engineering of the University of Porto (FEUP)}{Courses: 
% ``Linear Algebra'' (First Semester) and  ``Computing Systems'' (Second Semester)}{Feb. 2022 -- Ongoing}{Porto, Portugal}

% \cvevent{Student Monitor at the Federal Center for Technological Education of Rio de Janeiro (CEFET-RJ)}{Course: ``Physics III''}{2015}{Rio de Janeiro, Brasil}

% \begin{itemize}
% \item Student Travel Support offered by Coppetec by the work ``\textit{Dynamic Model and Line of Sight Control of a 3-DOF Inertial Stabilization Platform via Feedback Linearization}'', presented by me at the 2018 American Control Conference.
% \end{itemize}
% \cvevent{}{}{2018}{Milwaukee, USA}

%% If the NEXT page doesn't start with a \cvsection but you'd
%% still like to add a sidebar, then use this command on THIS
%% page to add it. The optional argument lets you pull up the
%% sidebar a bit so that it looks aligned with the top of the
%% main column.
%\addnextpagesidebar[-1ex]{CV-p2sidebar}